\documentclass[12pt]{article}
\usepackage{graphicx} % Required for inserting images
\usepackage[utf8]{inputenc} % Codificación de caracteres
\usepackage[spanish]{babel} % Idioma del documento
\usepackage{amsmath, amssymb, amsthm}
\usepackage{pgfplots}
\usepackage{tabularx}
\usepackage[margin=1in]{geometry}
\usepackage{svg}
\usepackage{float}
\usepackage{fancyhdr}
\usepackage{setspace}
\usepackage[most]{tcolorbox}

\newtcolorbox{comentarioprofe}[1]{
    colback=blue!5,
    colframe=blue!50!black,
    title=#1,
    fonttitle=\bfseries,
    boxrule=0.8pt,
    arc=2mm
}
\newtcolorbox{ejemplo}[1]{
    colback=gray!8,
    colframe=gray!70!black,
    title=#1,
    fonttitle=\bfseries,
    boxrule=0.8pt,
    arc=2mm
}
\newtcolorbox{importante}[1]{
    colback=orange!8,
    colframe=orange!70!black,
    title=#1,
    fonttitle=\bfseries,
    boxrule=0.8pt,
    arc=2mm
}

\geometry{margin=1in}
\onehalfspacing
\pagestyle{fancy}
\rhead{Macroeconomía II}
\lhead{Tema 2}
\title{Tema 2: Producción}
\author{Ignacio Sanabria}
\date{Agosto 2026}

\begin{document}
\maketitle

\section{Productividad}
En la parte de producción vamos a endogenizar $Y_t$\\
Suponga un modelo de dos periodos: 
$$Y_{t+j}=A_{t+j}F(K_{t+j},N_{t+j})\; \forall j=0,1$$
Donde: 
\begin{itemize}
    \item $K_{t+j}$= Stock físico
    \item $N_{t+j}$= Horas de trabajo 
    \item $F(.,.)$= Función de producción 
    \item $A_{t+j}$= \textbf{productividad total de factores}
\end{itemize}

\subsection{Capital Físico}
Sea $\delta>0$ una tasa de depreciación en un modelo de \textbf{2} periodos. \\
En el periodo $t$ hay una perdida de capital vista como: \\
$$K_t \xrightarrow{produccion} (1-\delta)K_t $$
Igual en el periodo futuro: \\
$$K_{t+1} \xrightarrow{produccion} (1-\delta)K_{t+1} $$
Vamos a asumir siempre que $K_t>0$, por el tiempo que lleva crear capital. \\
Si en el primer periodo tengo $K_t$, en el segundo periodo tendría: 
$$K_{t+1}=(1-\delta)K_t+I_t$$
Donde $I_t$ es el pedido de nuevo capital a futuro. También llamada \textbf{Ley de Movimiento del Capital}\\
\subsubsection{Ley de Movimiento del Capital}
\begin{align*}
    I_t=K_{t+1}-(1-\delta)K_t\\
    K_{t+2}=(1-\delta)+I_{t+1}
\end{align*}
Sin embargo por ser solo \textbf{2} periodos, quiero $K_{t+2}=0$, es decir: \\
$$ I_{t+1}=-(1-\delta)K_{t+1}$$
Por ende sabemos que $I_t$ es la \textbf{inversión en capital}, puede ser positiva si quiero obtener capital o negativa si quiero desecharlo. \\
De tal manera que: 
\begin{figure} [H]
        \centering 
        \includesvg[width=0.5\linewidth]{figuras/clase_2/macro_t2_1.drawio.svg}
        \label{fig:placeholder}
\end{figure}
\newpage
\subsection{Empresa}
Una empresa escoge $\{N_t,N_{t+1},I_t,I_{t+1},K_{t+1},K_{t+2}\}$ dado que: 
\begin{figure} [H]
        \centering 
        \includesvg[width=0.5\linewidth]{figuras/clase_2/macro_t2_2.drawio.svg}
        \label{fig:placeholder}
\end{figure}
Sea $D_t$= ganancias / dividendos que genera la empresa en $t$.
\subsubsection{Mercado Laboral}
\begin{itemize}
    \item La empresa demanda $N_{t+j} (N_{t+j}^d)$
    \item Hogares ofrecen $N_{t+j} (N_{t+j}^s)$
    \item Tranzan con salario real $w_{t+j}$, en base a bienes de consumo. 
\end{itemize}
\subsubsection{Creación de Capital}
La empresa tiene un $K_t$ dado, pero decide el $I_t$, sin embargo debe decidir si financiarlo con recursos propios o endeudamiento. \textbf{Es decir, existe un mercado de fondos prestables}\\
\subsubsection{Endeudarse}
Suponga que la empresa financia completamente su inversión en el mercado de fondos prestables. 
\begin{itemize}
    \item En $t$: Pide prestado $I_t$
    \item En $t+1$: Paga $I_t+\text{Interés} \quad (r_t^I)$\\
    Tal que $r_t^I=r_t+f_t$, donde $f_t$ es el margen de intermediación. 
\end{itemize}
\subsection{Utilidad}
Tomando todos estos supuestos sabemos que: 
$$D_t=Y_t-w_tN_t$$ son los dividendos en $t$\\
En $t+1$: \\
$$D_{t+1}=Y_{t+1}+(1+\delta)K_{t+1}-w_{t+1}N_{t+1}-(1-r_t^I)I_t$$
Por ende el problema de la empresa es: 
\begin{align*}
    \underset{\{N_t,N_{t+1},I_t,I_{t+1},K_{t+1},K_{t+2} \}}{max} \;\underbrace{D_t+\frac{D_{t+1}}{1+r_t}}_{\text{Valor de la empresa}}
\end{align*}
De manera mas sencilla se maximiza la siguiente expresión: 
$$\underbrace{\underbrace{A_tF(K_t,N_t)}_{Y_t}-w_tN_t}_{D_t}+\underbrace{(\frac{1}{1+r_t})[\underbrace{A_{t+1}F(K_{t+1},N_{t+1})}_{Y_{t+1}}+(1-\delta)K_{t+1}-w_{t+1}N_{t+1}-(1+r_t)\underbrace{(k_{t+1}-(1-\delta)K_t)}_{I_t}]}_{D_{t+1}}$$

Ahora: 
\begin{align*}
    [N_t]:& A_tF_N(K_t,N_t)-w_t=0\\
    [N_{t+1}]: & (\frac{1}{1+r_t})[A_{t+1}F_N(K_{t+1},N_{t+1})-w_{t+1}]=0 \\
    [K_{t+1}]: & (\frac{1}{1+r_t})[A_{t+1}F_N(K_{t+1},N_{t+1})+(1-\delta)-(1+r_t)=0
\end{align*}
\textbf{Condiciones de optimalidad}
\begin{enumerate}
    \item $A_tF_N(K_t,N_t)=w_t$
    \item $A_{t+1}F_N(K_{t+1},N_{t+1})= w_{t+1}$
    \item $A_{t+1}F_K(K_{t+1},N_{t+1})+(1-\delta)= (1+r_t^I)$
\end{enumerate}

\begin{comentarioprofe}{Observe:}
    Si $F(.,.)$ cumple las siguientes propiedades: 
    \begin{enumerate}
        \item $F(K,0)=0$ y $F(0,N)=0;\; \forall K,N$
        \item $F_k(K,N)>0$ y $F_N(K,N)>0;\; \forall K,N$
        \item $F_{KK}(K,N)<0$ y $F_{NN}(K,N)<0; \; \forall K,N$
        \item $F_{NK}(K,N)>0$ y $F_{KN}(K,N)>0; \; \forall K,N$
        \item $\forall\lambda : \lambda F(K,N)=F(\lambda K,\lambda N)$, es decir es homogénea de grado 1 y posee \textbf{rendimientos constantes a escala}
    \end{enumerate}
\end{comentarioprofe}
\begin{ejemplo}{Cobb Douglas}
    Sea: 
    $$Y_t=A_tK_T^\alpha N_t^{1-\alpha}$$
    Entonces: 
    \begin{itemize}
        \item COP.1: $A_t(1-\alpha)K_t^\alpha N^{-\alpha}_t=w_k$
        Es decir: $N_t=(\frac{(1-\alpha)A_tK_t}{w_k})^\frac{1}{\alpha}$
    \end{itemize}
    Es decir, sabemos que el comportamiento de $N_t$ se vera de forma: $$N_t=N^d(\underbrace{w_t}_{(-)},\underbrace{A_t}_{(+)},\underbrace{K_t}_{(+)})$$
    De igual forma: 
    \begin{itemize}
        \item COP.3 $\Rightarrow{}$ $K_{t+1}=k^d(\underbrace{r_t}_{(-)},\underbrace{A_{t+1}}_{(+)},\underbrace{f_t}_{(-)})$
        \item Recuerde que $I_t=k_{t+1}-(1-\delta)k_t$, de forma que: $$I_t=k^d(\underbrace{r_t}_{(-)},\underbrace{A_{t+1}}_{(+)},\underbrace{f_t}_{(-)})$$
    \end{itemize}
\end{ejemplo}
\newpage
\subsection{Hogar}
\begin{itemize}
    \item El hogar representativo es dueño de la empresa. $(D_t,D_{t+1})$, lo recibe de fir a exogena. 
    \item En terminos de una empresa, cuando produce mas que sus costos obtiene \textbf{dividendendos}, estos dividendos pertenecen a un hogar representativo.
    \item Entonces recibe: $(D_t,D_{t+1})$ de manera exógena. 
    \item El hogar obtiene utilidad del consumo y ocio. 
    \item El ocio se mide en horas. 
    \item Suponga un hogar que dispone de $H$ horas tal que: 
    $$H=\underbrace{L}_{Ocio}+\underbrace{N}_{Trabajo}$$
\end{itemize}
 Suponga que $u(C,L)$ con $u_c(C,L)>0;\; U_L(C,L)>0;\;U_{CC}(C,L)<0;\;U_{LL}(C,L)<0$\\
 Sujeto a: 
 \begin{itemize}
     \item En periodo $t:\quad C_t+S_t=w_tN_t+D_t$
     \item En periodo $t+1:\quad C_{t+1}=w_{t+1}N_{t+1}+D_{t+1}+(1+r_t)S_t$ donde $S_{t+1}=0$ por la condicion terminal. 
     \item Consolidando ambas restricciones: 
     $$C_t+\frac{C_{t+1}}{1+r_t}=w_tN_t+D_t+\frac{w_{t+1}N_{t+1}+D_{t+1}}{1+r_t}$$
 \end{itemize}
 Por ende el problema del hogar seria: 
 \begin{align*}
     max \;u(C_t,L_t)+\beta u(C_{t+1},L_{t+1})\; \\s.a\;\;\quad \quad \quad \quad \\ C_t+\frac{C_{t+1}}{1+r_t}=w_tN_t+D_t+\frac{w_{t+1}N_{t+1}+D_{t+1}}{1+r_t} \\
 \end{align*}
Como sabemos que $1=L+N$, entonces $L=1-N$\\
Por ende: 
 \begin{align*}
     max \;u(C_t,1-N_t)+\beta u(C_{t+1},1-N_{t+1})\; \\s.a\;\;\quad \quad \quad \quad \\ C_t+\frac{C_{t+1}}{1+r_t}=w_tN_t+D_t+\frac{w_{t+1}N_{t+1}+D_{t+1}}{1+r_t} \\
 \end{align*}
 Sabemos que el agente es menos feliz por trabajar ya que peirde ocio, sin embargo si no trabaja pierde capacidad de consumo, por lo cual debe encontrar un equilibrio de ambos. \\
 \begin{itemize}
     \item $[C_t]: U_c(C_t,1-N_t)-\lambda =0$\\
     \item $[C_{t+1}]: \beta U_c(C_{t+1}, 1 - N_{t+1})-\frac{\lambda}{1+r_t}=0$\\
     \item $[N_t]:-U_L(Ct,1-N_t)+\lambda w_t=0$
     \item $[\lambda]: RPI$
 \end{itemize}
 \begin{comentarioprofe}Observe
     Sea: 
     \begin{align*}
         u(C,L)=& log(C) + \theta \; log(L) \\
         u(C,1-N)=& log(C) + \theta \; log(1-N) \\
         \frac{\partial \;u(C,L)}{\partial \; N}=&\frac{u_L(C,L)}{\partial\; L} \frac{\partial\; L}{\partial\;N}=\frac{\partial(1-N)}{\partial N}=-1
     \end{align*}
     \textbf{¿Que mide $\theta$?}
 \end{comentarioprofe}

 Simplificando las CPO:\\
 \begin{enumerate}
    \item De [$C_t$] y [$C_{t+1}$]: 
     $U_c(C_t,1-N_t)=\beta(1+r_t)\;u_c(C_{t+1},1-N_{t+1})$
    \item De [$C_t$] y [$1-N_t$]: 
     $U_L(C_{t},1-N_{t})=w_t\;u_c(C_{t},1-N_{t})$
    \item De [$C_{t+1}$] y [$1-N_{t+1}$]: 
     $U_L(C_{t+1},1-N_{t+1})=w_{t+1}\;u_c(C_{t+1},1-N_{t+1})$
    \item RPI 
 \end{enumerate}

 Por lo tanto: 
 $$\underbrace{U_L(C_t,1-N_t)}_{\text{Beneficio marginal de una unidad de ocio}}=\underbrace{w_tU_C(C_t,1-N_t)}_{\text{Costo marginal de una unidad de ocio}}$$
 \textbf{Por ende $\theta$ mide la sustitucion entre $C_t$ y $N_t$}

\begin{ejemplo}Ejemplo 
    Suponga un modelo de un periodo sin ahorro. 
    $C_t=w_tN_t+D_t$\\
    Ademas $N_t=1-L_t$
    De forma que: 
    \begin{figure} [H]
            \centering 
            \includesvg[width=0.5\linewidth]{figuras/clase_2/macro_t2_5.drawio.svg}
            \label{fig:placeholder}
    \end{figure}
    Vamos a tener: $N_t=N^S(w_t,\theta_t) \Rightarrow{}$ \textbf{Funcion de oferta laboral}\\
    $$\frac{\partial\;N^S}{\partial\; w_t}\xrightarrow{CPO\;2},\;\frac{\partial\;N^S}{\partial\; w_t}>0 \Rightarrow{}\textbf{Efecto Sutitucion}\\\\
    \frac{\partial\;N^S}{\partial\; w_t}\xrightarrow{RPI},\;\uparrow C_t,\uparrow L_t\xrightarrow{}\downarrow N_t\Rightarrow{} \frac{\partial\;N^S}{\partial\; w_t}<0 \Rightarrow{}\textbf{Efecto Ingreso}
    $$
    \textbf{Vamos a asumir siempre que el efecto sustitucion domina.}\\
    $C_t=C^d(\underbrace{Y_t}_{(+)},\underbrace{Y_{t+1}}_{(+)},\underbrace{r_t}_{\underbrace{(-)}_{Efecto\; Sutitucion}})$
\end{ejemplo}
\end{document}